\chapter{The spectral namer as built}\label{ch:namer}

\begin{onething}
The spectral namer is a 347,685-parameter two-tower network on Indradhanu~V2's wavelength-aware stem: one tower reads
a pixel with its ring of neighbours, the other reads a laboratory spectrum, and both meet in one table of cosines;
since 2026-10-31 \emph{both} of its questions are answered from that table, the material class by pooling the
cosines of each man-made class, and natural-or-man-made by comparing the pooled man-made cosines with the pooled
natural ones, because a single linear head on the pixel tower sat at chance.
\end{onething}

\Cref{sec:material:twolevel} gave the \emph{why} of the redesigned naming step: use-based families, two questions
in turn, matching instead of classifying, and the separability check as a classical baseline. The design became
\file{docs/proposals/spectral-namer.md}, and then code, all on 2026-10-31. This chapter reads the code from end to
end, with every tensor shape and every parameter, and it tells the one change made after the first run started:
\emph{level~1 by matching}.

\paragraph{Status when this chapter was written.} Run N001 started on an H100 capacity block at 11:30~UTC on
2026-10-31 (block until 2026-11-02 11:30~UTC). After its first epoch its level-1 head was at chance; it was stopped
and replaced by N001b (\file{configs/namer/N001b.json}), the same run with level~1 by matching, and by N002, the
open-set run, also with matching. \textbf{The results of N001b and N002 are pending}: N001b's README records its
first epoch only, and N002 has none yet, so nothing in this chapter is a final result of either. Every number is a
setting \tagM{configs/namer/N001b.json}, a count or worked number from this book's script
\file{research/textbook/tools/compute_namer_full.py} \tagC, a CPU smoke test, or one of the early measurements now
recorded in \file{research/experiments/spectral-namer/N001/README.md} and \file{N001b/README.md}.

Paths abbreviate: \file{namer/} is \file{app/spectral_match/namer/}, and \file{components/} is
\file{app/foundation_models/components/}.

\begin{analogy}
A wine-tasting course that cannot find enough labelled wines. The teacher pours ordinary table wine, adds a measured
splash of a reference wine from the cellar, and asks two questions: ``was anything added that does not grow in a
vineyard?'' and ``which grape?''. The splash size varies, the glasses differ a little, the room is noisy. The student
is marked right for naming the \emph{grape}, not for recognising the exact bottle, and the bottle used for the splash
is never one of the answers on the card. The student answers both questions the same way: by asking which bottles on
the shelf the glass tastes most like. The table wine is a real pixel, the splash a library spectrum, the grape the
material class, and the shelf the library.
\end{analogy}

%======================================================================
\section{What the namer answers}\label{sec:namer:question}

The namer runs \emph{after} detection. Something else (\cref{ch:inference,ch:material}) has already flagged a pixel
as anomalous; the namer says what it might be. For each pixel it returns four numbers (\file{namer/infer.py},
\code{NamerResult}):
\begin{enumerate}
\item $P(M)$, the probability that something man-made is in the pixel: \emph{level~1};
\item the most likely man-made class $\hat k$ among seven;
\item $P(\hat k\mid M)$, that class's probability \emph{if} the pixel is man-made: \emph{level~2};
\item $\hat a$, an estimate of the \emph{fill}, the share of the pixel the material covers.
\end{enumerate}
The proposal combines the two levels as $P(k)=P(M)\times P(k\mid M)$. The seven classes are materials, not uses
(\file{namer/classes.py}, \code{MANMADE_CLASSES}): oil/fuel, plastic, textile/fabric, wood/paper, paint/coatings,
metal, concrete/asphalt.

\begin{definition}[Fill]
The \emph{fill}\index{fill} $a$ of a pixel is the fraction of its area covered by the target material, between 0 and
1. A 10\,m boat in a 30\,m pixel has $a\approx(10\times10)/(30\times30)=100/900=0.11$ \tagA.
\end{definition}

%======================================================================
\section{The architecture: a Siamese, two-tower network}\label{sec:namer:model}

\begin{definition}[Encoder, embedding, tower, Siamese network]
An \emph{encoder}\index{encoder} is a network that turns an input into a short vector, its
\emph{embedding}\index{embedding}. A \emph{tower}\index{tower} is one encoder in a model that has several, drawn
side by side. A \emph{Siamese network}\index{Siamese network} \citep{bromley1993signature} runs two inputs through
the \emph{same} weights and compares the outputs. A \emph{two-tower} (or dual-encoder\index{dual encoder}) model has
two towers whose outputs are compared, usually by a cosine; the towers may share some weights and not others.
\end{definition}

The namer is both, in a precise sense. Its first stage, the \emph{stem}, is Siamese: one set of 17,825 weights runs
on the pixel $\vect{x}$, on its ring mean $\vect{r}$ and on every library spectrum $\vect{t}_j$. After the stem the
two sides part into towers with different weights. The namer's own docstring (\file{namer/model.py}) says:
\begin{lstlisting}[language={}]
pixel side:   x, r (ring mean) -> stem -> [z_x, z_r, z_x - z_r] -> MLP -> h
              h -> embedding (unit, 128) | fill in [0, 1] | (level-1 linear head, unused with level1="matching")
library side: t -> stem -> MLP -> embedding (unit, 128)
class score:  logsumexp over the class's library embeddings of  cos(z, e_j) / tau
level 1:      "matching" (N001b on) - natural entries as one more class, level1_logit; or the linear head (N001)
\end{lstlisting}
The linear head is the one N001 used; N001b computes level~1 by matching (\cref{sec:namer:level1}), and the head is
still built but no longer read.

\begin{figure}[p]
\centering
\includegraphics[angle=90,height=0.86\textheight]{namer_network}
\caption{The publication schematic of the namer, from
\texttt{research/\allowbreak experiments/\allowbreak spectral-namer/\allowbreak figures/\allowbreak namer\_network.pdf}, redrawn on
31~October 2026 for level~1 by matching: the pixel tower has two heads (embedding and fill), and the class-matching
box pools the natural library entries as one more class, $P(\text{man-made})=\sigma(\operatorname{LSE}_{c<7}S_c-S_{\text{natural}}+b)$
(\cref{sec:namer:level1}, \cref{fig:namer:shapes}). The old Linear level-1 head is still built in the model but not used, so it is not drawn.}
\label{fig:namer:schematic}
\end{figure}

\subsection{The pixel tower}

The pixel tower encodes the pixel $\vect{x}$ and its ring mean $\vect{r}$ separately through the stem, getting
$\vect{z}_x$ and $\vect{z}_r$ of 64 numbers each, and concatenates
\begin{equation}
\vect{u}=[\vect{z}_x,\ \vect{z}_r,\ \vect{z}_x-\vect{z}_r]\in\R^{192}.
\label{eq:namer:concat}
\end{equation}
This is \cref{sec:material:naming}'s ``subtract the ring'', done on features rather than on reflectance; the network
also keeps $\vect{z}_x$ and $\vect{z}_r$ themselves, so it can learn when the subtraction helps and when it does not.
A three-layer perceptron (\code{_mlp}, Linear--GELU--Linear--GELU--Linear, \cref{ch:block} defined GELU) maps
$\vect{u}$ to $\vect{h}\in\R^{256}$. Three linear \emph{heads} read $\vect{h}$:
\begin{itemize}
\item \code{to_embed}: $256\to128$, then divided by its length, so the pixel embedding $\vect{z}$ is a unit vector;
\item \code{to_fill}: $256\to1$, then a sigmoid, so $\hat a\in(0,1)$;
\item \code{to_manmade}: $256\to1$, a raw score (a \emph{logit}). This was level~1 in N001; under
\code{"level1": "matching"} its output is overwritten before any loss sees it.
\end{itemize}

\subsection{The library tower}

The library tower runs each library spectrum $\vect{t}_j$ through the same stem, a second perceptron
$64\to256\to256\to128$ (\code{library}), and divides by the length: a unit vector $\vect{e}_j\in\R^{128}$.

\subsection{Where the towers meet}

Both embeddings have length 1, so their dot product is their cosine. With a learned temperature $\tau$:
\begin{equation}
s_j=\frac{\vect{z}\cdot\vect{e}_j}{\tau}=\frac{\cos(\vect{z},\vect{e}_j)}{\tau}.
\label{eq:namer:sj}
\end{equation}
For a batch of $P$ pixels and $L$ library entries, all the $s_j$ form one matrix
$\mat{S}=\mat{Z}\mat{E}\T/\tau$ of shape \shape{(P, L)}, the line \code{s = z @ lib_embed.T / self.log_tau.exp()} of
\code{class_scores}. Everything the namer says about materials is read from this matrix.

\subsection{Why the towers are not the same}

A Siamese network often uses one tower for both inputs. The namer does not, after the stem, because the two inputs
are different kinds of object:
\begin{itemize}
\item \textbf{The pixel is a mixture in context.} It is $a\vect{t}+(1-a)\vect{b}$ plus noise and an atmospheric
error, with $a$ often below 0.25. What identifies the target is the \emph{difference} from the surroundings, so the
pixel tower takes the ring too, and it must also estimate how much was added (the fill head).
\item \textbf{The library spectrum is pure and alone.} It is a laboratory measurement of one material, $a=1$, no
background, no ring, no noise worth modelling. The library tower has no ring input and no fill head.
\end{itemize}
The pixel tower's job is to \emph{undo} the mixing: to map a 15\%-plastic pixel over grass near the point where the
library tower maps pure plastic. The library tower only has to place pure spectra so that materials group together.
Forcing one set of weights to do both would make the library side pretend to un-mix something that was never mixed.
\textbf{The asymmetry is deliberate: the pixel tower learns to un-mix, the library tower only to arrange.}

\subsection{Why a new library needs no retraining}

Nothing in the model is tied to a library entry. The library tower is a function: spectrum in, embedding out. To add
a spectrum, or a whole class, you encode it and append the row to $\mat{E}$; to add a class you also give its rows a
new class index. Inference does exactly this every time: \code{name_pixels} builds the bank at the scene's own bands
and calls \code{forward_library} on all of it (\cref{sec:namer:infer}). A closed classifier, by contrast, has one
output per class; a new class means a new output and new training. The proposal's decision record says this is why
the namer matches rather than classifies (``a matching (not closed-classifier) model so a new library needs no
retraining'').

\begin{worked}[label={we:namer:newlib}]{The cost of a new library}
A new library of 500 spectra at EnMAP's bands, encoded on the laptop. The library tower is one stem pass over
\shape{(500, N, 1, 1)} and one perceptron over \shape{(500, 64)}. Its perceptron multiplies
\begin{align*}
64\times256+256\times256+256\times128 &= 16{,}384+65{,}536+32{,}768 = 114{,}688
\end{align*}
numbers per spectrum, so $500\times114{,}688=57{,}344{,}000$ multiply-adds: a fraction of a second on a CPU. Training
the namer is 20,000 steps of tens of thousands of samples. The difference is the reason for the design: a new library
costs one forward pass, not a training run \tagC.
\end{worked}

But note what ``no retraining'' does not promise. The library tower was trained on the family library's spectra; a new
spectrum that is unlike all of them may land somewhere unhelpful. Whether a new class is \emph{named well} without
training is exactly what the open-set run N002 measures (\cref{sec:namer:splits}).

\subsection{Counting the parameters}

\begin{worked}[label={we:namer:params}]{Counting the namer's parameters}
A linear layer from $i$ to $o$ numbers has $i\times o$ weights and $o$ biases. The pixel perceptron is
$192\to256\to256\to256$, the library perceptron $64\to256\to256\to128$:
\begin{align*}
\text{pixel MLP} &= (192\times256+256)+(256\times256+256)+(256\times256+256)\\
&= 49{,}408+65{,}792+65{,}792 = 180{,}992,\\
\text{library MLP} &= (64\times256+256)+(256\times256+256)+(256\times128+128)\\
&= 16{,}640+65{,}792+32{,}896 = 115{,}328,\\
\text{heads} &= \underbrace{(256\times128+128)}_{\text{embedding}}+\underbrace{2\times(256+1)}_{\text{manmade, fill}}
 = 32{,}896+514 = 33{,}410,\\
\text{BatchNorm} &= \underbrace{64}_{\text{scale}}+\underbrace{64}_{\text{shift}} = 128,\\
\text{scalars} &= \underbrace{1}_{\ln\tau}+\underbrace{1}_{\text{level-1 bias }b} = 2,\\
\text{stem} &= 17{,}825 \quad\text{(counted by hand in \cref{sec:namer:stemparams})},\\
\text{total} &= 180{,}992+115{,}328+33{,}410+128+2+17{,}825 = 347{,}685 .
\end{align*}
Instantiating \code{SpectralNamer} with N001b's settings gives 347,685 \tagC. The stem, the only part taken from
Indradhanu, is $17{,}825/347{,}685=5.1\%$ of it. The pixel tower (perceptron and three heads) is
$180{,}992+33{,}410=214{,}402$, the library tower 115,328.
\end{worked}

\begin{sourcenote}[: 347,685, not 347,684; and not 1--3\,M]
The first version of this chapter, and the course, counted 347,684. That was right for the code before level~1 by
matching: \code{l1_bias}, one number, was added on 2026-10-31 with \code{level1_logit}. The proposal's ``Size: small
($\sim$1-3 M parameters)'' is three to nine times the model in the code. Of the 347,685, the 257 of
\code{to_manmade} receive no gradient under matching (\cref{sec:namer:level1}).
\end{sourcenote}

%======================================================================
\section{The stem in full, at the namer's shapes}\label{sec:namer:stem}

\Cref{sec:specdecoder:reduced} built the reduced stem at $D=32$ for SpectralEarth. The namer runs it at
$D=64$, Indradhanu Heavy's width, copied from the V2 checkpoint FTH3 epoch~5 by \code{from_spectral_stem} with the
settings \code{compressed_channels} 64, \code{token_dim} 32, \code{attention_heads} 4, \code{bin_nm} 40,
\code{band_weighting} \code{"width"} (FTH3's own config has the same five values
\tagM{research/experiments/indradhanu-finetune/FTH3-heavy-joint-more2/config.json}). This section walks it once more,
at the shapes it really runs in one namer training step: $P=8192$ pixels (or $L=378$ library spectra) and $N=200$
bands \tagA.

The namer calls it through \code{SpectralNamer.encode}:
\begin{lstlisting}
x = spectra[:, :, None, None]                     # (P, N) -> (P, N, 1, 1): P images of one pixel
z = self.stats.normalize(x, wavelengths_nm)       # per-wavelength z-score
...
return self.norm(self.stem(z, wavelengths_nm, bv)[:, :, 0, 0])   # (P, 64)
\end{lstlisting}
The stem works along wavelength inside each pixel, so a batch of $P$ spectra is passed as $P$ images of size
$1\times1$. Nothing spatial happens.

\subsection{Step 0: normalisation by wavelength}

\code{WavelengthStats} (\file{components/pixel_normalization.py}) holds a mean curve $\mu(\lambda)$ and a spread curve
$\sigma(\lambda)$ over wavelength, measured once on SpectralEarth's training split and stored in the checkpoint
(``never recomputed''). A band at wavelength $\lambda$ between two stored wavelengths $\lambda_{lo}<\lambda<\lambda_{hi}$
reads the curves by linear interpolation:
\begin{align}
t &= \frac{\lambda-\lambda_{lo}}{\lambda_{hi}-\lambda_{lo}}, \qquad
\mu(\lambda)=\mu_{lo}+t\,(\mu_{hi}-\mu_{lo}), \qquad \sigma(\lambda)=\sigma_{lo}+t\,(\sigma_{hi}-\sigma_{lo}),\\
v &= \frac{x-\mu(\lambda)}{\sigma(\lambda)} .
\end{align}
Outside the stored range the end values are held.

\begin{worked}[label={we:namer:lerp}]{Normalising one band}
Stored: at 860\,nm mean 0.30, spread 0.10; at 870\,nm mean 0.32, spread 0.12 \tagA. A band at 864\,nm reads 0.35:
\begin{align*}
t &= (864-860)/(870-860) = 4/10 = 0.4,\\
\mu &= 0.30+0.4\times(0.32-0.30) = 0.30+0.008 = 0.308,\\
\sigma &= 0.10+0.4\times(0.12-0.10) = 0.10+0.008 = 0.108,\\
v &= (0.35-0.308)/0.108 = 0.042/0.108 = 0.389 .
\end{align*}
\end{worked}

This keeps \emph{absolute} reflectance meaningful: a bright and a dark pixel stay different, unlike a per-scene
min-max scaling. The proposal names this as a requirement (``it keeps absolute reflectance''). Invalid bands are set
to 0 after normalisation (\code{masked_fill}) and are also excluded from every average below.

\subsection{Step 1: the wavelength encoding}

Every band's wavelength becomes a vector of $d=32$ numbers (\file{components/spectral_tokens.py},
\code{wavelength_encoding}): 16 periods spaced geometrically from 10 to 10,000\,nm, a sine and a cosine of each,
\begin{equation}
p_i=10\times1000^{\,i/15}\ \text{nm},\quad i=0,\dots,15;\qquad
\operatorname{enc}(\lambda)=\Big[\sin\tfrac{2\pi\lambda}{p_0},\dots,\sin\tfrac{2\pi\lambda}{p_{15}},\
\cos\tfrac{2\pi\lambda}{p_0},\dots,\cos\tfrac{2\pi\lambda}{p_{15}}\Big].
\label{eq:namer:enc}
\end{equation}
Successive periods differ by the factor $1000^{1/15}=1.5849$: 10.0, 15.8, 25.1, 39.8, \dots, 6309.6, 10000.0\,nm \tagC.

\begin{worked}[label={we:namer:enc}]{Three entries of $\operatorname{enc}(1000\text{ nm})$}
\begin{align*}
\sin(2\pi\times1000/10) &= \sin(200\pi) = 0,\\
\sin(2\pi\times1000/15.85) &= \sin(396.4) = 0.5659,\\
\sin(2\pi\times1000/10000) &= \sin(0.6283) = 0.5878,\qquad \cos(0.6283)=0.8090 .
\end{align*}
The code gives the same \tagC. The slowest pair (period 10,000\,nm) changes by only a few per cent between
neighbouring bands, so it says roughly \emph{where} in the spectrum a band sits; the fastest pair separates bands 5\,nm
apart. The encodings of 1000 and 1010\,nm have cosine 0.6709 \tagC: similar, not identical.
\end{worked}

For the batch, $\operatorname{enc}$ of all $N$ bands is one matrix of shape \shape{(N, 32)} = \shape{(200, 32)}, the
same for every pixel; it is computed once per call.

\subsection{Step 2: bins of 40\,nm}

The bins are a fixed grid from 400\,nm, the same on every sensor (\code{BIN_ORIGIN_NM}, \code{BIN_END_NM} = 400, 2500):
\begin{equation}
G=\Big\lceil\frac{2500-400}{40}\Big\rceil=\lceil52.5\rceil=53,\qquad
g(\lambda)=\operatorname{clamp}\Big(\Big\lfloor\frac{\lambda-400}{40}\Big\rfloor,\ 0,\ 52\Big).
\end{equation}
So 405 and 439.9\,nm fall in bin 0, 440\,nm in bin 1, 1000\,nm in bin $\lfloor600/40\rfloor=15$, 2450\,nm in bin 51; the
last bin runs from 2480 to 2520\,nm, past the grid's end, and a band below 400\,nm is clamped into bin 0 \tagC. Bin $g$
has centre $420+40g$\,nm. The membership is a one-hot matrix $\mat{O}$ of shape \shape{(N, G)} = \shape{(200, 53)}:
$O_{bg}=1$ if band $b$ is in bin $g$.

On EnMAP's real 224-band list (418--2445\,nm), 47 of the 53 bins hold at least one band, at most 9 in one bin; the
empty ones are the two water-vapour gaps and the bins past 2445\,nm \tagC.

\subsection{Step 3: band widths}

With \code{band_weighting="width"} each band counts by the nanometres of spectrum it stands for (\code{band_widths}):
half the gap to each neighbour, the single gap at either end, every gap capped at \code{MAX_GAP_NM} = 20\,nm.

\begin{worked}[label={we:namer:widths}]{Widths of four bands, and of a band beside a water gap}
Bands at 410, 430, 445, 470\,nm. Gaps: $430-410=20$, $445-430=15$, $470-445=25$, capped to 20, 15, 20.
\begin{align*}
w_{410} &= 20 &&\text{one gap at the end}\\
w_{430} &= (20+15)/2 = 17.5,\qquad w_{445} = (15+20)/2 = 17.5\\
w_{470} &= 20 &&\text{one (capped) gap at the end}
\end{align*}
Bands at 1320, 1330, 1340, 1460, 1470\,nm: the 120\,nm water gap is capped to 20, so the widths are 10, 10,
$(10+20)/2=15$, $(20+10)/2=15$, 10. Without the cap the 1340\,nm band would claim 65\,nm of spectrum it never measured.
The code gives both lists \tagC.
\end{worked}

The widths, times validity, form the weight $c_b$ of band $b$ for each pixel, and the stem multiplies the values by it:
\code{vals} $=c_b v_b$, shape \shape{(P, N)}; \code{cnt_band} $=c_b$, shape \shape{(P, N)}; and the total weight per
bin \code{cnt} $=\mat{C}_{\text{band}}\mat{O}$, shape \shape{(P, N)}$\times$\shape{(N, G)} = \shape{(P, 53)}.

\subsection{Step 4: bin tokens, as one matrix product}

Each band also gets its own weight vector $\vect{k}(\lambda_b)\in\R^{32}$, from a small network of its encoding
(\code{band_weight}, $32\to64\to32$). Bin $g$'s token is the weighted mean of its bands' $v_b\,\vect{k}(\lambda_b)$:
\begin{equation}
\vect{\tau}_g=\frac{\sum_{b\in g}c_b\,v_b\,\vect{k}(\lambda_b)}{\sum_{b\in g}c_b}.
\label{eq:namer:token}
\end{equation}
Written as a loop this would build a \shape{(P, N, 32)} tensor, 52.4 million numbers at $P=8192$, $N=200$. The code
avoids it by folding the bin membership into the weights. Build
\begin{equation}
\mat{K}\in\R^{N\times(G\cdot d)},\qquad K_{b,\,(g,i)}=O_{bg}\,k_i(\lambda_b),
\end{equation}
row $b$ holding $\vect{k}(\lambda_b)$ in the 32 columns of its own bin and zeros elsewhere
(\code{kmat = (onehot[:, :, None] * self.band_weight(enc)[:, None, :]).reshape(n, G * self.d)}). Then one product
\shape{(P, N)}\,@\,\shape{(N, 53$\cdot$32)} = \shape{(P, 1696)}, reshaped to \shape{(P, 53, 32)} and divided by
\code{cnt}, gives every token of every pixel. $\mat{K}$ has $200\times1696=339{,}200$ entries; the tokens
$8192\times53\times32=13{,}893{,}632$ \tagC.

\begin{worked}[label={we:namer:kmat}]{Tokens by the matrix product, on paper}
Take the four bands of \cref{we:namer:widths} with $d=2$ instead of 32 and two bins (410, 430 in bin~0; 445, 470 in
bin~1). Normalised values $v=(0.5,-0.2,1.0,0.4)$ and band vectors $\vect{k}=(1,0)$, $(0.5,0.5)$, $(0,1)$, $(1,-1)$
\tagA. Widths $c=(20,17.5,17.5,20)$.
\begin{align*}
c\,v &= (20\times0.5,\ 17.5\times(-0.2),\ 17.5\times1.0,\ 20\times0.4) = (10,\ -3.5,\ 17.5,\ 8),\\
\text{cnt} &= (20+17.5,\ 17.5+20) = (37.5,\ 37.5),\\
\mat{K} &= \begin{pmatrix}1&0&0&0\\0.5&0.5&0&0\\0&0&0&1\\0&0&1&-1\end{pmatrix},\\
(c\,v)\,\mat{K} &= \big(10\times1+(-3.5)\times0.5,\ \ (-3.5)\times0.5,\\
&\qquad\ 8\times1,\ \ 17.5\times1+8\times(-1)\big)\\
&= (8.25,\ -1.75,\ 8.0,\ 9.5),\\
\vect{\tau}_0 &= (8.25,\ -1.75)/37.5 = (0.2200,\ -0.0467),\\
\vect{\tau}_1 &= (8.0,\ 9.5)/37.5 = (0.2133,\ 0.2533).
\end{align*}
A direct loop over the bands of each bin gives the same two tokens \tagC.
\end{worked}

Then every token gets its bin's position,
\begin{lstlisting}
tok = tok + self.bin_bias(wavelength_encoding(centres))     # (P, 53, 32) + (53, 32)
\end{lstlisting}
a linear map $32\to32$ of the encoding of the bin centre, shape \shape{(53, 32)}, added to every pixel's tokens.

\subsection{Step 5: 64 queries attend over 53 bins, in 4 heads}

The stem has 64 learned queries, a \shape{(64, 32)} matrix, one per output channel. Attention (\cref{ch:attention})
lets each query read the bins it needs:
\begin{center}\small
\begin{tabular}{@{}lll@{}}
\toprule
what & formula & shape\\
\midrule
queries & $\mat{Q}=\operatorname{Lin}_q(\text{queries})$, all pixels & \shape{(64, 32)} $\to$ 4 heads \shape{(64, 8)}\\
keys, values & $\mat{K}_p=\operatorname{Lin}_k(\vect{\tau})$, $\mat{V}_p=\operatorname{Lin}_v(\vect{\tau})$ & \shape{(P, 53, 32)} $\to$ \shape{(P, 4, 53, 8)}\\
weights & $\softmax\big(\mat{Q}\mat{K}_p\T/\sqrt8\big)$, masked & \shape{(P, 4, 64, 53)}\\
answers & weights $\times\,\mat{V}_p$, heads joined & \shape{(P, 64, 32)}\\
\bottomrule
\end{tabular}
\end{center}
The mask is \code{cnt > 0}: a bin with no valid band gets weight exactly 0 (\cref{sec:crossstem:mask}). If a pixel
has no valid band at all, the mask is dropped for it and its output is set to 0 at the end.

\begin{worked}[label={we:namer:attn}]{One query, three bins, one empty}
Head dimension 2 for paper: query $\vect{q}=(1,0.5)$; keys $(1,0)$, $(0,1)$, $(2,2)$; values $(0.2,0.4)$, $(0.6,-0.2)$,
$(9,9)$; bin~3 has no valid band \tagA.
\begin{align*}
\vect{q}\cdot\vect{k}/\sqrt2 &= (1,\ 0.5,\ 3)/1.4142 = (0.7071,\ 0.3536,\ 2.1213),\\
\text{masked} &= (0.7071,\ 0.3536,\ -\infty),\\
e^{0.7071} &= 2.0281,\quad e^{0.3536}=1.4242,\quad e^{-\infty}=0,\quad\text{sum }3.4523,\\
\text{weights} &= (2.0281/3.4523,\ 1.4242/3.4523,\ 0) = (0.5875,\ 0.4125,\ 0),\\
\text{answer} &= 0.5875\,(0.2,0.4)+0.4125\,(0.6,-0.2) = (0.3650,\ 0.1525).
\end{align*}
The large value 9 in the empty bin cannot leak. PyTorch's \code{scaled_dot_product_attention} with the same boolean
mask returns $(0.3650,\ 0.1525)$ \tagC.
\end{worked}

If the attention weights were stored, they would be $8192\times4\times64\times53=111{,}149{,}056$ numbers, 424\,MiB in
float32, per stem call \tagC. PyTorch's fused attention kernel avoids storing them, and the stem splits calls above
32,768 rows (\code{SDPA_MAX_ROWS}); at $P=8192$ one call suffices.

\subsection{Step 6: one number per channel, plus the skip path}

The answers pass through $\operatorname{Lin}_o$ ($32\to32$) and \code{to_channel} ($32\to1$): \shape{(P, 64, 32)}
$\to$ \shape{(P, 64, 1)} $\to$ \shape{(P, 64)}. Channel $c$ is query $c$'s answer squeezed to one number.

Alongside, a per-band \emph{skip path} lets every valid band reach the channels directly with weights computed from
its wavelength (\code{skip_weight}, $32\to64\to64$):
\begin{equation}
\text{skip}_c=\frac{\sum_b c_b\,v_b\,W_c(\lambda_b)}{\sum_b c_b},
\end{equation}
in code \code{(vals @ self.skip_weight(enc)) / cnt_band.sum(1)}: \shape{(P, N)} @ \shape{(N, 64)} = \shape{(P, 64)}.
This is the old $1\times1$ compressor with columns computed from wavelength; attention adds adaptivity on top, it does
not carry all the detail alone (the stem's docstring, ``fix 2''). The stem returns
$\vect{z}=\text{attention}+\text{skip}$, \shape{(P, 64)} (as \shape{(P, 64, 1, 1)}).

\subsection{Step 7: the frozen BatchNorm}

\begin{definition}[BatchNorm in evaluation mode]
A BatchNorm\index{BatchNorm} layer \citep{ioffe2015batch} with stored running mean $m_c$ and variance $s_c^2$, learned
scale $\gamma_c$ and shift $\beta_c$, maps channel $c$ of an input as
$y_c=\gamma_c\,(z_c-m_c)/\sqrt{s_c^2+\epsilon}+\beta_c$, $\epsilon=10^{-5}$. In \emph{training} mode it uses the
current batch's mean and variance instead of $m_c,s_c^2$, and updates the stored ones.
\end{definition}

\begin{worked}[label={we:namer:bn}]{One channel through the frozen BatchNorm}
$m=0.4$, $s^2=0.25$, $\gamma=1.5$, $\beta=-0.1$, input $z=0.9$ \tagA:
\begin{align*}
z-m &= 0.9-0.4 = 0.5,\qquad \sqrt{0.25+0.00001} = 0.50001,\\
y &= 1.5\times0.5/0.50001-0.1 = 1.49997-0.1 = 1.39997 .
\end{align*}
\code{BatchNorm1d} in evaluation mode gives 1.39997 \tagC.
\end{worked}

\code{SpectralNamer.train()} calls the usual \code{train()} and then \code{self.norm.eval()}: the running statistics
stay those of V2, estimated on real pixels of the three sensors, and normalise pixels, rings and library spectra alike,
in training and evaluation; $\gamma$ and $\beta$ still train. \Cref{sec:namer:loss} tells why.

\subsection{The stem's parameters, by hand}\label{sec:namer:stemparams}

\begin{table}[htbp]
\centering\small
\begin{tabular}{llr}
\toprule
piece & layers & parameters\\
\midrule
\code{band_weight}, $\vect{k}(\lambda)$ & $32\to64$, $64\to32$ & $(32\times64+64)+(64\times32+32)=2{,}112+2{,}080=4{,}192$\\
\code{skip_weight}, $W(\lambda)$ & $32\to64$, $64\to64$ & $(32\times64+64)+(64\times64+64)=2{,}112+4{,}160=6{,}272$\\
\code{bin_bias} & $32\to32$ & $32\times32+32=1{,}056$\\
\code{queries} & $64\times32$ & $2{,}048$\\
$\operatorname{Lin}_{q,k,v,o}$ & $4\times(32\to32)$ & $4\times1{,}056=4{,}224$\\
\code{to_channel} & $32\to1$ & $33$\\
\midrule
total & & $4{,}192+6{,}272+1{,}056+2{,}048+4{,}224+33=17{,}825$\\
\bottomrule
\end{tabular}
\caption{The reduced stem at $D=64$, $d=32$. Instantiation gives 17,825, tensor by tensor \tagC; the Indradhanu
course's part~71 derives the same total. No count depends on the number of bands $N$ or on the sensor.}
\label{tab:namer:stemparams}
\end{table}

\begin{table}[htbp]
\centering\small
\begin{tabular}{lll}
\toprule
step & tensor & shape\\
\midrule
input & spectra as images & \shape{(8192, 200, 1, 1)}\\
0 & normalised values $v$, masked & \shape{(8192, 200, 1, 1)}\\
1 & $\operatorname{enc}(\lambda)$ & \shape{(200, 32)}\\
2 & one-hot bins $\mat{O}$ & \shape{(200, 53)}\\
3 & \code{vals}, \code{cnt_band}; \code{cnt} & \shape{(8192, 200)}; \shape{(8192, 53)}\\
4 & $\vect{k}(\lambda)$; $\mat{K}$; tokens & \shape{(200, 32)}; \shape{(200, 1696)}; \shape{(8192, 53, 32)}\\
4 & bin bias & \shape{(53, 32)}\\
5 & $\mat{Q}$; $\mat{K}_p,\mat{V}_p$ & \shape{(1, 4, 64, 8)} expanded; \shape{(8192, 4, 53, 8)}\\
5 & weights (not stored); answers & \shape{(8192, 4, 64, 53)}; \shape{(8192, 64, 32)}\\
6 & \code{to_channel}; skip; $\vect{z}$ & \shape{(8192, 64, 1)}; \shape{(8192, 64)}; \shape{(8192, 64)}\\
7 & BatchNorm (frozen) & \shape{(8192, 64)}\\
\bottomrule
\end{tabular}
\caption{The stem for the 8,192 pixels of one sensor's batch, $N=200$ bands \tagA. Forward hooks on the real modules
record these shapes \tagC; for the library replace 8192 by 378.}
\label{tab:namer:stemshapes}
\end{table}

%======================================================================
\section{The ring mean}\label{sec:namer:ring}
\index{ring mean}

A pixel's surroundings are summarised by the mean of the valid pixels in a $15\times15$ box minus the central
$5\times5$ box (\file{namer/generator.py}, \code{INNER, OUTER = 5, 15}). Box sums come for every pixel at once from an
average pool times the box area:
\begin{lstlisting}
s = box(v, OUTER) - box(v, INNER)     # sum of valid values in the ring, (B, N, H, W)
c = box(m, OUTER) - box(m, INNER)     # number of valid pixels in the ring
ring = s / c.clamp(min=1)
\end{lstlisting}
In the interior $c=15^2-5^2=225-25=200$. The $5\times5$ box covers offsets $-2$ to $2$, so it removes every pixel up to
two steps away (a diagonal step counts as one: the \emph{Chebyshev distance}\index{Chebyshev distance}
$\max(|\Delta\text{row}|,|\Delta\text{col}|)$). The ring is the pixels at Chebyshev distance 3 to 7. Check: the square
ring at distance $d$ holds $(2d+1)^2-(2d-1)^2=8d$ pixels, and $8\times(3+4+5+6+7)=8\times25=200$. The two-pixel gap keeps
a target's halo (\cref{ch:benchmarks}) out of its own background.

\begin{worked}[label={we:namer:ring}]{A ring that ignores the halo}
A flat tile of 0.20; the target pixel 0.50; its eight neighbours 0.35 \tagA. The $15\times15$ box holds the
target, its 8 neighbours and $225-9=216$ background pixels; the $5\times5$ box holds the same 9 and $25-9=16$
background pixels:
\begin{align*}
S_{15} &= 216\times0.20+8\times0.35+0.50 = 43.2+2.8+0.5 = 46.5,\\
S_{5} &= 16\times0.20+8\times0.35+0.50 = 3.2+2.8+0.5 = 6.5,\\
\vect{r} &= \frac{S_{15}-S_5}{200} = \frac{40.0}{200} = 0.200 .
\end{align*}
The target and its halo cancel exactly; the code gives 0.2000 with a count of 200 \tagC.
\end{worked}

Samples are drawn only where every band is valid, the ring has at least 20 valid pixels and the pixel is at least
7~px from the tile's edge, so of a $128\times128$ tile at most $(128-14)^2=12{,}996$ pixels qualify \tagC. Inference
computes the same ring one band at a time with \code{scipy}'s \code{uniform_filter} (\file{namer/infer.py},
\code{ring_means}); a test pins the two to $10^{-5}$.

\begin{sourcenote}[: 3 to 7 pixels, not 2 to 7]
The proposal's diagram says ``r (ring mean 2-7 px)''. The code's ring is Chebyshev distance 3 to 7, as the
generator's docstring and, now, \file{infer.py}'s docstring say (``3-7 px''); an earlier version of \file{infer.py}
said ``2-7 px''.
\end{sourcenote}

%======================================================================
\section{The generator}\label{sec:namer:generator}
\index{planted-sample generator}

Almost no real pixel comes with a known material, so the namer trains on samples it makes from real tiles of the
three sensors' fine-tuning shards. For a sampled pixel $\vect{b}$ with ring mean $\vect{r}$ (\file{namer/generator.py},
\code{plant}):
\begin{equation}
\vect{x}=\vect{g}\odot\big(a\,\vect{t}+(1-a)\,\vect{b}\big)+\vect{n},\qquad \vect{r}'=\vect{g}\odot\vect{r},
\label{eq:namer:plant}
\end{equation}
where $\odot$ multiplies band by band. The kind of sample decides $\vect{t}$ and $a$:
\begin{center}\small
\begin{tabular}{llll}
\toprule
kind & share & $\vect{t}$ & level-1 truth\\
\midrule
0, man-made & 0.50 & a man-made library spectrum of the training split & man-made\\
1, natural mix & 0.25 & a natural library spectrum (either split) & natural\\
2, plain & 0.25 & none, $a=0$ & natural\\
\bottomrule
\end{tabular}
\end{center}
(\code{kinds} = 0.5, 0.25, 0.25 \tagM{configs/namer/N001b.json}.) The kind is drawn per sample with
\code{torch.multinomial}, so the counts vary: one draw of 8,192 gave 4,110, 1,984 and 2,098 \tagC. Each sample records
its library entry index (\code{entry}, $-1$ for plain), its class (\code{cls}, $-1$ unless man-made) and its fill.

\paragraph{Picking pixels.} For each tile the generator draws 512 positions (\code{pixels_per_tile}) uniformly,
\emph{with replacement}, among the allowed pixels of \cref{sec:namer:ring}; with 16 tiles that is $P=16\times512=8192$
samples per sensor per step. It gathers $\vect{b}$ and $\vect{r}$ at those positions, shape \shape{(8192, N)} each,
and the tile's noise level $\sigma$, one per band, also \shape{(8192, N)}.

\paragraph{The fill, log-uniform.} With $u$ uniform on $[0,1]$, the code sets
$a=\exp\big(\ln a_{\min}\,(1-u)\big)=a_{\min}^{\,1-u}$ with $a_{\min}=0.05$: $u=0$ gives $a_{\min}$, $u=1$ gives 1.
Its distribution, one step per line; remember that $\ln a_{\min}<0$:
\begin{align*}
a\le q\ &\Longleftrightarrow\ (1-u)\ln a_{\min}\le\ln q &&\text{take logs}\\
&\Longleftrightarrow\ 1-u\ \ge\ \frac{\ln q}{\ln a_{\min}} &&\text{divide by $\ln a_{\min}<0$: it flips}\\
&\Longleftrightarrow\ u\ \le\ 1-\frac{\ln q}{\ln a_{\min}} &&\text{rearrange}\\
P(a\le q) &= 1-\frac{\ln q}{\ln a_{\min}} &&\text{$P(u\le v)=v$, $u$ uniform}\\
&= \frac{\ln a_{\min}-\ln q}{\ln a_{\min}} &&\text{common denominator}\\
&= \frac{\ln q-\ln a_{\min}}{-\ln a_{\min}} &&\text{times }\tfrac{-1}{-1}
\end{align*}
which, with $\ln q-\ln a_{\min}=\ln(q/a_{\min})$ and $-\ln a_{\min}=\ln(1/a_{\min})$, is
\begin{equation}
P(a\le q)=\frac{\ln(q/a_{\min})}{\ln(1/a_{\min})} .
\label{eq:namer:loguniform}
\end{equation}
A shorter way to see it: $\ln a=(1-u)\ln a_{\min}$ moves linearly from $\ln a_{\min}$ to 0 as $u$ goes from 0 to 1,
so $\ln a$ is uniform on $[\ln a_{\min},0]$, and the share below $\ln q$ is that interval's length fraction. That is
why the distribution is called \emph{log-uniform}\index{log-uniform distribution}.

\paragraph{Its density and mean.} The density is the slope of \cref{eq:namer:loguniform} in $q$. The slope of
$\ln q$ is $1/q$, so
\begin{equation}
p(a)=\frac{1}{a\,\ln(1/a_{\min})}\quad\text{for }a_{\min}\le a\le1 .
\end{equation}
The mean, one step per line:
\begin{align*}
\E[a] &= \int_{a_{\min}}^{1} a\,p(a)\,da &&\text{definition of a mean}\\
&= \int_{a_{\min}}^{1} a\cdot\frac{1}{a\ln(1/a_{\min})}\,da &&\text{insert the density}\\
&= \frac{1}{\ln(1/a_{\min})}\int_{a_{\min}}^{1} 1\,da &&\text{the $a$ cancels}\\
&= \frac{1-a_{\min}}{\ln(1/a_{\min})} &&\text{the integral of 1 is the interval's length.}
\end{align*}

\begin{worked}[label={we:namer:fill}]{How much training goes to small fills}
With $a_{\min}=0.05$, $\ln(1/0.05)=\ln20=2.9957$. By \cref{eq:namer:loguniform}:
\begin{align*}
P(a\le0.1) &= \frac{\ln(0.1/0.05)}{2.9957}=\frac{\ln2}{2.9957}=\frac{0.6931}{2.9957}=0.2314,\\
P(a\le0.25) &= \frac{\ln5}{2.9957}=\frac{1.6094}{2.9957}=0.5372,\\
P(a\le0.5) &= \frac{\ln10}{2.9957}=\frac{2.3026}{2.9957}=0.7686 .
\end{align*}
The mean: $\E[a]=(1-0.05)/2.9957=0.95/2.9957=0.3171$; the median is where $P=\frac12$, $u=\frac12$, so
$a=0.05^{0.5}=0.2236$. A million samples drawn with the code's formula put 0.5368 of them at or below 0.25, median
0.2241 \tagC. A uniform fill on $[0.05,1]$ would put only $(0.25-0.05)/0.95=0.21$ there and have mean 0.525.
\Cref{tab:material:separability} showed small fills are where ACE fails; the log-uniform spends more than half of the
man-made training there.
\end{worked}

The fill \emph{target} of the fill head is $a$ for kinds 0 and 1 and 0 for plain pixels, so over all samples its mean
is $0.75\times0.3171=0.2378$ \tagC.

\paragraph{The gain, smooth.} Three coefficients $c_0,c_1,c_2\sim\mathcal{N}(0,0.5^2)$ make a parabola over the band
index rescaled to $t\in[-1,1]$, clipped to $[-1,1]$; then $g=1+0.1\times\text{curve}$, so $0.9\le g\le1.1$
(\code{smooth_gain}, shape \shape{(P, N)}). It stands for an atmospheric-correction error, and the same $\vect{g}$
multiplies the ring mean, because the atmosphere errs alike over a pixel and its neighbours. With $c=(0.2,-0.4,0.6)$
\tagA: at $t=-1$ the curve is $0.2+0.4+0.6=1.2$, clipped to 1, $g=1.10$; at $t=0$, $g=1.02$; at $t=1$,
$0.2-0.4+0.6=0.4$, $g=1.04$.

\paragraph{The noise, from the tile.} Per tile and band, $\sigma$ is the root mean square of the differences between
horizontal neighbours, divided by $\sqrt2$ (\code{ring_and_noise}, shape \shape{(B, N)}). The reason, using one rule:
\emph{the variance of a difference of independent variables is the sum of their variances.} Check it on two coins:
$X,Y\in\{0,1\}$ each with probability $\frac12$ have variance $\frac14$; $X-Y$ takes $-1,0,0,1$ with equal chance, mean
0, variance $\frac14(1+0+0+1)=\frac12=\frac14+\frac14$. So if neighbours carry independent noise of variance
$\sigma^2$,
\begin{align*}
\operatorname{Var}(n_1-n_2) &= \sigma^2+\sigma^2 = 2\sigma^2 &&\text{the rule}\\
\operatorname{sd}(n_1-n_2) &= \sigma\sqrt2 &&\text{square root}\\
\sigma &= \operatorname{sd}(n_1-n_2)/\sqrt2 &&\text{divide by }\sqrt2 .
\end{align*}

\begin{worked}[label={we:namer:noise}]{Noise from four pixels}
One band of a row reads $0.20, 0.23, 0.21, 0.22$ \tagA:
\begin{align*}
\text{differences} &= 0.03,\ -0.02,\ 0.01\\
\text{squares} &= 0.0009,\ 0.0004,\ 0.0001, \qquad \text{sum } 0.0014\\
\text{mean square} &= 0.0014/3 = 0.000467\\
\text{RMS} &= \sqrt{0.000467} = 0.0216\\
\sigma &= 0.0216/1.4142 = 0.01528 .
\end{align*}
The code gives 0.01528; on 200,000 samples of white noise with $\sigma=0.01$ it estimates 0.01003 \tagC.
\end{worked}

\begin{sourcenote}[: an RMS, not a standard deviation]
The differences' mean is not subtracted before squaring (``RMS (uncentred) \dots\ texture counts too'', the
docstring), so real texture, an edge or a slope, counts as noise. On textured tiles $\sigma$ is an over-estimate, which
errs on the hard side.
\end{sourcenote}

\begin{worked}[label={we:namer:plant}]{Planting one pixel}
Three bands, $\vect{b}=(0.10,0.30,0.25)$, $\vect{t}=(0.40,0.35,0.05)$, $a=0.25$, $\vect{g}=(1.02,1.00,0.98)$,
$\vect{n}=(0.003,-0.002,0.001)$ \tagA. By \cref{eq:namer:plant}:
\begin{align*}
a\vect{t} &= (0.1000,\ 0.0875,\ 0.0125),\\
(1-a)\vect{b} &= (0.0750,\ 0.2250,\ 0.1875),\\
\text{sum} &= (0.1750,\ 0.3125,\ 0.2000),\\
\vect{g}\odot\text{sum} &= (1.02\times0.1750,\ 1.00\times0.3125,\ 0.98\times0.2000) = (0.1785,\ 0.3125,\ 0.1960),\\
\vect{x} &= (0.1785+0.003,\ 0.3125-0.002,\ 0.1960+0.001) = (0.1815,\ 0.3105,\ 0.1970) .
\end{align*}
\end{worked}

What is \emph{not} modelled matters as much: the planting is one pixel, with no point-spread function, no halo and no
adjacency, and the ring is untouched by the target. The proposal's ``Risks'' list says so (``linear mixing, smooth gain
and Gaussian noise are simplifications''). Per batch only bands valid in more than half the tiles' pixels are kept
(\code{keep = valid.float().mean(dim=(0, 2, 3)) > 0.5}), so water-vapour bands drop out, and the library is cut to the
same bands (\code{bank.bands(keep)}); a batch with fewer than 20 bands left is skipped.

\begin{sourcenote}[: no band-list jitter]
The proposal lists ``Band-list jitter as in pretraining (\texttt{band\_augment\_prob})''. Neither the generator nor the
trainer reads such a setting; every sample of a batch has the batch's live bands.
\end{sourcenote}

%======================================================================
\section{The class score}\label{sec:namer:classscore}

\begin{definition}[Logsumexp, class score]
For numbers $z_1,\dots,z_m$, $\operatorname{LSE}(z)=\ln\sum_j e^{z_j}$\index{logsumexp}, a smooth maximum: it lies
between $\max_j z_j$ and $\max_j z_j+\ln m$. With $s_j$ from \cref{eq:namer:sj}, the \emph{class score} of class $k$ is
$S_k=\operatorname{LSE}(\{s_j: j\in k\})$, and $P(k\mid M)$ is the softmax of $S_1,\dots,S_7$ (\code{class_scores}).
\end{definition}

The temperature\index{temperature} $\tau$ is learned, stored as $\ln\tau$ and started at $\ln0.07=-2.6593$, so
$1/\tau=14.286$ at the start: a cosine of 0.9 against 0.7 becomes a score of 12.86 against 10.00.

\paragraph{The class softmax is a sum of entry softmaxes.} Write $q_j=e^{s_j}/\sum_i e^{s_i}$ for the softmax over all
library entries that belong to some class. Then
\begin{align*}
P(k\mid M) &= \frac{e^{S_k}}{\sum_{k'}e^{S_{k'}}} &&\text{definition}\\
&= \frac{\sum_{j\in k}e^{s_j}}{\sum_{k'}\sum_{j\in k'}e^{s_j}} &&e^{\ln u}=u\\
&= \frac{\sum_{j\in k}e^{s_j}}{\sum_i e^{s_i}} &&\text{every entry is in exactly one class}\\
&= \sum_{j\in k}q_j .
\end{align*}
So the loss treats \emph{any} member of the right class as a right answer. That is the remedy the companion volume
proposed for ``false negatives'' in a per-spectrum contrastive loss, two spectra of one material pushed apart; here it
is built in.

\paragraph{Leaving the planted spectrum out.} In \code{forward}, \code{class_scores(..., exclude=s["entry"])} sets the
planted entry's $s_j$ to $-\infty$ before the logsumexp, in training and in evaluation alike. Every planted sample
(man-made and natural mix) has its own entry excluded; plain pixels have none. A test checks that a class whose only
member is excluded scores $-\infty$.

\begin{worked}[label={we:namer:classscore}]{Recognising the sample, or naming the material}
Take $\tau=0.1$ for round numbers \tagA. The pixel was planted with metal spectrum $e_1$; cosines: metal $e_1$ 0.95,
$e_2$ 0.70, $e_3$ 0.60; plastic $e_4$ 0.75; oil $e_5$ 0.20. Divided by $\tau$: 9.5, 7.0, 6.0, 7.5, 2.0.

With $e_1$ in the library:
\begin{align*}
S_{\text{metal}} &= \ln(e^{9.5}+e^{7}+e^{6}) = 9.6064, \quad S_{\text{plastic}}=7.5, \quad S_{\text{oil}}=2.0,\\
P &= \operatorname{softmax}(9.6064,\ 7.5,\ 2.0) = (0.8911,\ 0.1084,\ 0.0004),\\
L &= -\ln0.8911=0.1153 .
\end{align*}
With $e_1$ left out:
\begin{align*}
S_{\text{metal}} &= \ln(e^{7}+e^{6}) = 7+\ln(1+e^{-1}) = 7+\ln1.3679 = 7.3133,\\
P &= (0.4524,\ 0.5453,\ 0.0022),\qquad L=-\ln0.4524=0.7931 .
\end{align*}
Check the sum-of-entries identity: over $e_2,\dots,e_5$ the entry softmax is $(0.3308,0.1217,0.5453,0.0022)$, and
$0.3308+0.1217=0.4525$, the metal probability up to rounding \tagC.
\end{worked}

With its own spectrum present the pixel is ``metal'' at 0.89 because it recognises $e_1$. With $e_1$ out it must lean
on the other metals, and here plastic wins. \textbf{Leaving the planted spectrum out makes the only route to a low loss
putting all metals near each other: learning the material, not the sample.} It is the separability check's
leave-one-spectrum-out (\code{ace[e] = -1}) used as a training rule.

\paragraph{The gradient.}\label{par:namer:l2grad} Let $c$ be the true class. The loss is
$L=-S_c+\ln\sum_{k}e^{S_k}=-\ln\sum_{j\in c}e^{s_j}+\ln\sum_i e^{s_i}$. Differentiate with respect to one entry's
$s_j$:
\begin{align*}
\frac{\partial}{\partial s_j}\ln\sum_i e^{s_i} &= \frac{e^{s_j}}{\sum_i e^{s_i}} = q_j &&\text{chain rule}\\
\frac{\partial}{\partial s_j}\Big(-\ln\sum_{i\in c}e^{s_i}\Big) &= -\,\mathbf{1}[j\in c]\,w_j,\quad
 w_j=\frac{e^{s_j}}{\sum_{i\in c}e^{s_i}} &&\text{the same rule, inside class }c\\
\frac{\partial L}{\partial s_j} &= q_j-\mathbf{1}[j\in c]\,w_j .
\end{align*}
Wrong-class entries are pushed down by their probability $q_j$, as in the per-spectrum loss. The true class's members
are pulled up, but shared out by $w_j$: the member already closest gets most of the pull.

\begin{worked}[label={we:namer:gradient}]{Who gets the pull}
In \cref{we:namer:classscore} with $e_1$ out: within metal, $w=(e^7,e^6)/(e^7+e^6)=(0.7311,\ 0.2689)$. With
$q=(0.3308,0.1217,0.5453,0.0022)$ for $e_2,e_3,e_4,e_5$:
\begin{align*}
e_2 &: 0.3308-0.7311=-0.4003, & e_3 &: 0.1217-0.2689=-0.1473,\\
e_4 &: 0.5453, & e_5 &: 0.0022 .
\end{align*}
They sum to $-0.4003-0.1473+0.5453+0.0022=-0.0001\approx0$, as they must \tagC. Plastic $e_4$, the tempting wrong
answer, is pushed hardest; metal's pull is shared $73:27$.
\end{worked}

Which entries take part in level~2? In training, only \emph{training-split man-made} entries (\code{train_lib}, minus an
open-set class). Test-split entries are still encoded (the whole bank goes through the library tower) but get class
index $-1$, which no class selects.

\begin{sourcenote}[: class-level, not per-spectrum InfoNCE]
The proposal's ``Losses'' (InfoNCE, \citealp{oord2018representation}) says ``InfoNCE over the batch's library spectra.
The positive is the planted spectrum''. The code does what its ``As built'' list says instead: cross-entropy over class
scores, own spectrum left out. Read the proposal's loss section as superseded.
\end{sourcenote}

%======================================================================
\section{Level 1 by matching}\label{sec:namer:level1}

\subsection{What N001 did, and what it measured}

In N001, level~1 was the linear head \code{to_manmade}: one weight vector $\vect{w}\in\R^{256}$ and a bias on the
pixel tower's $\vect{h}$, $\ell=\vect{w}\cdot\vect{h}+b_0$, $P(M)=\sigma(\ell)$, trained by binary cross-entropy. Two
measurements, both from 2026-10-31 \tagM{research/experiments/spectral-namer/N001/README.md}:
\begin{itemize}
\item \textbf{N001, epoch 1, on the H100.} Level-1 accuracy was about 0.50 at every evaluation fill on all three
sensors: on EnMAP 0.52 at fill 0.25 and 0.50 at fill 1.0. In the same epoch level-2 class accuracy was 0.60 to 0.68.
\item \textbf{A laptop CPU test, 150 steps at full fill on the EnMAP smoke shards}, the two level-1 forms side by side:
\end{itemize}
\begin{table}[htbp]
\centering\small
\begin{tabular}{lccc}
\toprule
level 1 & BCE & accuracy & man-made recall at ACE's false-alarm rate\\
\midrule
chance & $\ln2=0.6931$ & 0.50 & --\\
linear head on $\vect{h}$ & 0.64 & 0.65 & 0.41\\
by matching & 0.43 & 0.80 & 0.74\\
\bottomrule
\end{tabular}
\caption{The laptop test that decided the switch, as recorded in N001's README \tagM{research/experiments/spectral-namer/N001/README.md}. ACE at fill 1.0 on
EnMAP calls 13.36\% of natural plantings man-made and recalls 0.8711 of man-made ones \tagC, so after 150 steps neither
form is yet at ACE's level; the comparison is between the two forms.}
\label{tab:namer:l1evidence}
\end{table}

Read the epoch-1 numbers together. At fill 1.0 a planted man-made pixel \emph{is} a library spectrum (with gain and
noise), and level~2 named its class correctly about two times in three: the embedding $\vect{z}$, computed from the
same $\vect{h}$, plainly carried the information. The linear head on $\vect{h}$ could not read ``man-made'' out of it.
\textbf{The information was there; one linear layer was the wrong way to read it.}

To compare the two BCEs, turn each into a probability: $e^{-\text{BCE}}$ is the geometric mean of the probability given
to the truth. Chance: $e^{-0.6931}=0.5000$; linear: $e^{-0.64}=0.5273$; matching: $e^{-0.43}=0.6505$ \tagC.

N001 was stopped. N001b is the same configuration plus \code{"level1": "matching"} (\file{configs/namer/N001b.json},
whose description says ``N001's linear level-1 head sat at chance after epoch 1''), and N002 adds it too.

\subsection{The rule}

Under \code{"level1": "matching"} (\file{namer/train.py}, \code{forward}) the natural library entries become one more
class, index $C=7$:
\begin{lstlisting}
lib_cls = torch.where(bank.manmade, lib_cls, torch.full_like(lib_cls, C))     # natural -> class 7
scores = self.model.class_scores(z.float(), e.float(), lib_cls, C + 1, exclude=s["entry"])   # (P, 8)
logit, scores = self.model.level1_logit(scores), scores[:, :C]
\end{lstlisting}
and \code{level1_logit} (\file{namer/model.py}) is
\begin{lstlisting}
return torch.logsumexp(scores[:, :-1], dim=1) - scores[:, -1] + self.l1_bias
\end{lstlisting}
In symbols, with $S_0,\dots,S_6$ the man-made class scores, $S_7=S_{\text{nat}}$ the natural one, and $b$ the learned
scalar \code{l1_bias} (started at 0):
\begin{equation}
\ell=\operatorname{LSE}(S_0,\dots,S_6)-S_{\text{nat}}+b,\qquad P(M)=\sigma(\ell)=\frac{1}{1+e^{-\ell}} .
\label{eq:namer:l1}
\end{equation}
The level-2 scores are the first seven columns, exactly as before; the old head's output is discarded.

\subsection{What the rule means}

Two identities make \cref{eq:namer:l1} readable.

\emph{A logsumexp of logsumexps is the logsumexp of everything.} Since $S_c=\ln\sum_{j\in c}e^{s_j}$,
\begin{align*}
\operatorname{LSE}(S_0,\dots,S_6) &= \ln\sum_{c<7}e^{S_c} &&\text{definition}\\
&= \ln\sum_{c<7}\sum_{j\in c}e^{s_j} &&e^{\ln u}=u\\
&= \ln\sum_{j\in\mathcal{M}}e^{s_j} &&\mathcal{M}\text{: all man-made entries in play.}
\end{align*}
Check on five numbers $(1,2,0.5,3,-1)$ in groups $\{1,2\},\{0.5\},\{3,-1\}$: both ways give 3.4723 \tagC.

\emph{The sigmoid of a log-ratio is a share.} Write $A=\sum_{j\in\mathcal{M}}e^{s_j}$ and
$N=\sum_{j\in\mathcal{N}}e^{s_j}$ over the natural entries. Then $\ell=\ln A-\ln N+b$, and
\begin{align*}
e^{\ell} &= e^{b}A/N &&\text{exponentiate}\\
\sigma(\ell) &= \frac{e^{\ell}}{e^{\ell}+1} &&\text{multiply top and bottom of }\tfrac{1}{1+e^{-\ell}}\text{ by }e^{\ell}\\
&= \frac{e^{b}A/N}{e^{b}A/N+1} = \frac{e^{b}A}{e^{b}A+N} &&\text{multiply top and bottom by }N .
\end{align*}
So
\begin{equation}
P(M)=\frac{e^{b}\sum_{j\in\mathcal{M}}e^{s_j}}{e^{b}\sum_{j\in\mathcal{M}}e^{s_j}+\sum_{j\in\mathcal{N}}e^{s_j}} :
\label{eq:namer:l1share}
\end{equation}
\textbf{$P(M)$ is the share of the pixel's softmax mass over the library that falls on man-made entries}, with the
man-made side weighted by $e^{b}$. Level~1 asks ``does this pixel look more like the man-made shelf or the natural
shelf?'', in the same embedding level~2 uses.

\subsection{Its gradient, through the logsumexp}

The level-1 loss is $L_1=\operatorname{BCE}(\ell,y)=-[y\ln\sigma(\ell)+(1-y)\ln(1-\sigma(\ell))]$ with $y=1$ for
man-made. Two pieces, then the chain rule.

\emph{The BCE's slope in the logit.} With $p=\sigma(\ell)$ and $\sigma'(\ell)=p(1-p)$:
\begin{align*}
\frac{\partial L_1}{\partial\ell} &= -\Big[y\,\frac{p(1-p)}{p}-(1-y)\,\frac{p(1-p)}{1-p}\Big] &&\text{slope of }\ln\text{ is }1/\cdot\\
&= -[y(1-p)-(1-y)p] &&\text{cancel}\\
&= -[y-yp-p+yp] = p-y &&\text{expand.}
\end{align*}

\emph{The logit's slope in each entry score.} For $j\in\mathcal{M}$, only $\ln A$ depends on $s_j$; for
$j\in\mathcal{N}$, only $-\ln N$:
\begin{align*}
\frac{\partial\ell}{\partial s_j} &= \frac{e^{s_j}}{A}=:u_j &&j\in\mathcal{M}\ \text{(softmax within the man-made entries)}\\
\frac{\partial\ell}{\partial s_j} &= -\frac{e^{s_j}}{N}=:-v_j &&j\in\mathcal{N}\ \text{(softmax within the natural entries)}\\
\frac{\partial\ell}{\partial b} &= 1 .
\end{align*}

\emph{Together}, by the chain rule:
\begin{equation}
\frac{\partial L_1}{\partial s_j}=(p-y)\,u_j\ \ (j\in\mathcal{M}),\qquad
\frac{\partial L_1}{\partial s_j}=-(p-y)\,v_j\ \ (j\in\mathcal{N}),\qquad
\frac{\partial L_1}{\partial b}=p-y .
\label{eq:namer:l1grad}
\end{equation}
Since $\sum_{\mathcal{M}}u_j=1$ and $\sum_{\mathcal{N}}v_j=1$, the entry gradients sum to $(p-y)-(p-y)=0$: like level~2's,
a step moves score from one side to the other, it cannot raise everything. Finally
$s_j=\vect{z}\cdot\vect{e}_j/\tau$, so $\partial s_j/\partial\vect{z}=\vect{e}_j/\tau$ and the pixel embedding moves
along $\sum_j(\partial L_1/\partial s_j)\,\vect{e}_j/\tau$: \emph{towards} the man-made entries it is already closest to
and \emph{away} from the nearest natural ones when it is man-made, and the reverse when it is natural.

\begin{worked}[label={we:namer:l1}]{Level 1 by matching, three library entries}
One pixel, truth man-made ($y=1$). Three entries in play: plastic $s=6.0$, metal $s=5.0$ (both man-made, in different
classes) and soil $s=5.5$ (natural); $b=0$ \tagA.

The logit:
\begin{align*}
S_{\text{plastic}} &= 6.0,\quad S_{\text{metal}}=5.0,\quad S_{\text{nat}}=5.5 &&\text{one member each}\\
\operatorname{LSE}(6.0,5.0) &= 6+\ln(1+e^{-1}) = 6+\ln1.3679 = 6+0.3133 = 6.3133\\
\ell &= 6.3133-5.5+0 = 0.8133 .
\end{align*}
The probability, two ways:
\begin{align*}
p=\sigma(0.8133) &= 1/(1+e^{-0.8133}) = 1/(1+0.4434) = 0.6928,\\
\text{share} &= \frac{e^{6}+e^{5}}{e^{6}+e^{5}+e^{5.5}} = \frac{403.4288+148.4132}{403.4288+148.4132+244.6919}
= \frac{551.8420}{796.5339} = 0.6928 .
\end{align*}
The loss: $L_1=-\ln0.6928=0.3670$. The gradient, by \cref{eq:namer:l1grad}, with $p-y=0.6928-1=-0.3072$:
\begin{align*}
u &= \big(e^{6},e^{5}\big)/(e^{6}+e^{5}) = (403.4288,\ 148.4132)/551.8420 = (0.7311,\ 0.2689),\\
\partial L_1/\partial s_{\text{plastic}} &= -0.3072\times0.7311 = -0.2246,\\
\partial L_1/\partial s_{\text{metal}} &= -0.3072\times0.2689 = -0.0826,\\
\partial L_1/\partial s_{\text{soil}} &= -(-0.3072)\times1 = +0.3072,\\
\partial L_1/\partial b &= -0.3072 .
\end{align*}
Sum: $-0.2246-0.0826+0.3072=0$. Autograd on \code{level1_logit} gives $(-0.2246,\ -0.0826,\ 0.3072)$ \tagC. Gradient
descent raises the plastic score most (the closer man-made entry), the metal score a little, lowers the soil score, and
raises $b$.

Had the truth been natural ($y=0$): $L_1=-\ln(1-0.6928)=1.1803$ and the gradients flip sign and scale by $p=0.6928$:
$(0.5065,\ 0.1863,\ -0.6928)$ \tagC.
\end{worked}

\subsection{Why matching helps}

A hypothesis, consistent with \cref{tab:namer:l1evidence} but not proven by it. Three reasons.

\paragraph{1. One hyperplane must hold all of nature on one side.} A linear head says ``man-made'' when
$\vect{w}\cdot\vect{h}+b_0>0$: one flat cut through the 256-dimensional space. Nature is not one cluster: vegetation,
soils, rocks, water, snow and dry plant matter sit in different places, and man-made materials sit between them (a
green tarpaulin near vegetation, concrete near rock). For a \emph{fixed} feature space, one cut can be impossible.

\begin{worked}[label={we:namer:xor}]{Four points no line can split}
Two dimensions, unit vectors: natural at $0^\circ$ and $180^\circ$, $\vect{n}_1=(1,0)$, $\vect{n}_2=(-1,0)$; man-made at
$90^\circ$ and $270^\circ$, $\vect{m}_1=(0,1)$, $\vect{m}_2=(0,-1)$ \tagA. Suppose a line $\vect{w}\cdot\vect{h}+b_0$ is
negative on both natural points and positive on both man-made ones:
\begin{align*}
w_1+b_0<0,\quad -w_1+b_0<0 &\ \Rightarrow\ \text{adding: } 2b_0<0,\\
w_2+b_0>0,\quad -w_2+b_0>0 &\ \Rightarrow\ \text{adding: } 2b_0>0 .
\end{align*}
Both cannot hold: no line does it. Matching does, with the four points as library entries and $\tau=0.1$. A pixel at
$80^\circ$ has cosines $0.9848$ to $\vect{m}_1$, $-0.9848$ to $\vect{m}_2$, $0.1736$ to $\vect{n}_1$, $-0.1736$ to
$\vect{n}_2$:
\begin{align*}
\operatorname{LSE}(9.848,\ -9.848) &= 9.848,\qquad \operatorname{LSE}(1.736,\ -1.736)=1.767,\\
\ell &= 9.848-1.767 = 8.081\ \Rightarrow\ \text{man-made}.
\end{align*}
A pixel at $190^\circ$ gets $\ell=-8.081$, natural \tagC. The nearest shelf decides.
\end{worked}

The perceptron beneath the head \emph{could} bend $\vect{h}$ until one cut works; that is what training would have
to discover from one bit per sample. Matching needs no such bend: it is a nearest-prototype rule
\citep{snell2017prototypical}, and the namer has 219 natural prototypes on EnMAP's bank \tagC, one per natural
spectrum, rather than one half-space.

\paragraph{2. It reuses what level 2 already learns.} The class loss shapes $\vect{z}$ and $\vect{e}_j$ with a rich
signal: seven classes, own spectrum left out. A separate linear head on $\vect{h}$ competes with that for the same
features. Under matching, level~1 reads the \emph{same} cosine table, so whatever arrangement helps level~2 helps level~1;
and level~1's gradient (\cref{eq:namer:l1grad}) is the same kind of signal, moving the pixel toward and away from
specific entries. The natural entries now also get a training signal they never had: in N001 nothing pulled a pixel
toward a natural spectrum.

\paragraph{3. Nature gets described by examples.} A new natural spectrum in the library is a new prototype at
inference without retraining, as for man-made ones. A linear head's idea of nature is frozen in $\vect{w}$.

\paragraph{One built-in offset.} The logsumexp of $m$ equal scores is $s+\ln m$, so the side with more entries scores
higher when nothing distinguishes them. On EnMAP's bank the training uses $219$ natural and $159-34=125$ man-made
entries: an offset of $\ln(219/125)=\ln219-\ln125=5.3891-4.8283=0.5608$ towards ``natural'' \tagC. The learned bias $b$
can absorb it; it is one reason $b$ exists.

\begin{sourcenote}[: what changes in the code and docs]
\code{to_manmade} is still built and still computed by \code{forward_pixel}; under matching its output is replaced
before the loss, so its 257 parameters receive no gradient (with \code{zero_grad(set_to_none=True)} AdamW skips them,
weight decay included). \code{load_namer} reads \code{config["level1"]} from the checkpoint (default \code{"head"}) and
accepts a checkpoint without \code{l1_bias}, so N001's checkpoints still load. The model's docstring now describes
both forms, and the proposal's ``As built'' list records level~1 by matching; the proposal's design text (``Level 1 does
not use the library at all'') and the publication schematic (\cref{fig:namer:schematic}) still show the old head. The
code and \file{configs/namer/N001b.json} are what runs.
\end{sourcenote}

\paragraph{What level 1 by matching does not change.} Natural spectra are planted in both splits, so level-1 accuracy
on natural plantings is not an unseen-spectrum test; but the planted natural entry is left out of $S_{\text{nat}}$
(\code{exclude}), so a natural mix cannot be recognised as its own spectrum either. And the planted man-made test
spectra at evaluation are matched only against \emph{training-split} man-made entries, because the class scores read
the training library (\cref{sec:namer:pass}).

%======================================================================
\section{One training step, end to end}\label{sec:namer:step}

This section is where the pieces stop being separate. One step of \code{NamerRun.train} loops over the three sensors;
for each it draws one batch of tiles, plants, encodes, scores and back-propagates its loss divided by 3. Here is one
sensor's share at illustrative EnMAP sizes: $B=16$ tiles of $128\times128$, 512 pixels per tile so $P=8192$
\tagM{configs/namer/N001b.json}; $N=200$ live bands \tagA\ (EnMAP's list has 224; a batch keeps those valid in more than
half its pixels, and the true count depends on the tiles); $L=378$ library entries \tagC; $C=7$; $D=64$; $d=32$; 4
heads; $G=53$. The shapes in \cref{tab:namer:step1,tab:namer:step2} were recorded by forward hooks on the real code with
random data at these sizes \tagC.

\begin{table}[htbp]
\centering\small
\begin{tabular}{@{}lp{0.42\linewidth}l@{}}
\toprule
\# & operation (file, function) & output shape\\
\midrule
1 & decode tiles: int16 $\times10^{-4}$, $-32768$ invalid (\code{decode}) & \shape{(16, 224, 128, 128)}\\
2 & keep live bands; bank to the same bands (\code{forward}) & \shape{(16, 200, 128, 128)}; bank \shape{(378, 200)}\\
3 & box sums: ring mean, ring count (\code{ring_and_noise}) & \shape{(16, 200, 128, 128)}, \shape{(16, 1, 128, 128)}\\
4 & noise $\sigma$ from horizontal differences & \shape{(16, 200)}\\
5 & allowed pixels: all bands valid, $\ge20$ ring pixels, 7 from the edge & \shape{(16, 128, 128)} bool\\
6 & pick 512 per tile with replacement; gather $\vect{b}$, $\vect{r}$, $\sigma$ & \shape{(8192, 200)} each\\
7 & kind per sample (0.5, 0.25, 0.25); entry; fill & \shape{(8192,)} each\\
8 & $\vect{t}$ = bank rows (zeros for plain); gain $\vect{g}$; noise & \shape{(8192, 200)} each\\
9 & $\vect{x}=\vect{g}\odot(a\vect{t}+(1-a)\vect{b})+\vect{n}$, $\vect{r}'=\vect{g}\odot\vect{r}$ & \shape{(8192, 200)} each\\
\bottomrule
\end{tabular}
\caption{One step, part 1: the generator, for one sensor. The tile tensor alone is $16\times200\times128\times128
=52{,}428{,}800$ numbers, 200\,MiB in float32 \tagC.}
\label{tab:namer:step1}
\end{table}

\begin{table}[htbp]
\centering\small
\begin{tabular}{@{}lp{0.42\linewidth}l@{}}
\toprule
\# & operation & output shape\\
\midrule
10 & stem + frozen BN on $\vect{x}$ (\cref{tab:namer:stemshapes}) & $\vect{z}_x$: \shape{(8192, 64)}\\
11 & stem + frozen BN on $\vect{r}'$ & $\vect{z}_r$: \shape{(8192, 64)}\\
12 & concatenate $[\vect{z}_x,\vect{z}_r,\vect{z}_x-\vect{z}_r]$ & \shape{(8192, 192)}\\
13 & pixel MLP & $\vect{h}$: \shape{(8192, 256)}\\
14 & \code{to_embed}, unit length; \code{to_fill}, sigmoid & $\vect{z}$: \shape{(8192, 128)}; $\hat a$: \shape{(8192,)}\\
15 & stem + frozen BN on all 378 library spectra & \shape{(378, 64)}\\
16 & library MLP, unit length & $\mat{E}$: \shape{(378, 128)}\\
17 & $\mat{S}=\mat{Z}\mat{E}\T/\tau$ & \shape{(8192, 378)}, 3,096,576 numbers\\
18 & own entry $\to-\infty$ (every planted sample) & \shape{(8192, 378)}\\
19 & class index per entry: train man-made 0--6, natural 7, test $-1$ & \shape{(378,)}\\
20 & pool: logsumexp over each class's columns & \shape{(8192, 8)}\\
21 & level 1: $\operatorname{LSE}$(cols 0--6) $-$ col 7 $+\,b$ & $\ell$: \shape{(8192,)}\\
22 & level 2: cols 0--6 & \shape{(8192, 7)}\\
23 & BCE($\ell$, man-made), all samples & scalar\\
24 & CE(level 2, class), man-made samples only ($\approx4096$ rows) & scalar\\
25 & $0.5\times$ mean $|\hat a-a|$, all samples & scalar\\
26 & $(L_1+L_2+0.5L_f)/3$, backward & --\\
\bottomrule
\end{tabular}
\caption{One step, part 2: the model and the losses, for one sensor. The stem runs three times per sensor (rows 10, 11,
15), nine times per step. After all three sensors: clip the gradient norm to 1, one AdamW step, one scheduler step.}
\label{tab:namer:step2}
\end{table}

\begin{figure}[htbp]
\centering
\begin{tikzpicture}[
  font=\footnotesize, >={Stealth[length=2mm]},
  b/.style={draw, rounded corners=2pt, align=center, inner sep=3pt, minimum height=7mm},
  gen/.style={b, fill=gray!8, draw=maskgrey},
  st/.style={b, fill=spec!12, draw=spec},
  px/.style={b, fill=spat!8, draw=spat},
  lb/.style={b, fill=green!8, draw=green!45!black},
  mt/.style={b, fill=attn!10, draw=attn},
  ls/.style={b, fill=boxwarn, draw=loss, text=loss!70!black},
  sh/.style={font=\scriptsize\ttfamily, text=gray!70!black}]
% generator
\node[gen, text width=60mm] (tiles) at (0,0) {16 tiles, live bands\\ \texttt{(16, 200, 128, 128)}};
\node[gen, text width=60mm, below=4mm of tiles] (plant) {ring mean, noise $\sigma$, pick 512 per tile, plant\\
$\vect{x}=\vect{g}\odot(a\vect{t}+(1-a)\vect{b})+\vect{n}$,\quad $\vect{r}'=\vect{g}\odot\vect{r}$};
\node[lb, text width=36mm] (bank) at (6.2,0) {library bank\\ \texttt{(378, 200)}};
\draw[->] (tiles) -- (plant);
\draw[->, dashed] (bank.south west) to[bend left=10] node[sh, above, sloped] {$\vect{t}$ rows} (plant.north east);
% stems
\node[st, text width=26mm, below left=6mm and -26mm of plant] (sx) {stem + BN on $\vect{x}$\\ \texttt{(8192, 64)}};
\node[st, text width=26mm, right=4mm of sx] (sr) {stem + BN on $\vect{r}'$\\ \texttt{(8192, 64)}};
\node[st, text width=36mm] (st) at (6.2,-2.65) {stem + BN on $\vect{t}_j$\\ \texttt{(378, 64)}};
\draw[->] (plant.south -| sx.north) -- node[sh, left] {\texttt{(8192, 200)}} (sx);
\draw[->] (plant.south -| sr.north) -- (sr);
\draw[->] (bank) -- (st);
\draw[<->, spec, dashed] (sr.east) -- node[sh, above] {shared} (st.west);
% towers
\node[px, text width=60mm, below=5mm of sx.south west, anchor=north west] (cat)
  {$[\vect{z}_x,\vect{z}_r,\vect{z}_x-\vect{z}_r]$ \texttt{(8192, 192)} $\to$ pixel MLP $\to\vect{h}$ \texttt{(8192, 256)}};
\node[px, text width=24mm, below=4mm of cat.south west, anchor=north west] (fill) {$\hat a$, sigmoid\ \texttt{(8192,)}};
\node[px, text width=28mm, right=3mm of fill] (emb) {$\vect{z}$, unit\ \texttt{(8192, 128)}};
\node[b, draw=maskgrey, dashed, text=maskgrey, text width=18mm, right=4mm of emb] (old)
  {\texttt{to\_manmade}\ not used};
\draw[->, maskgrey, dashed] (cat.east) -| (old.north);
\node[lb, text width=36mm] (lm) at (6.2,-4.6) {library MLP, unit\\ $\mat{E}$ \texttt{(378, 128)}};
\draw[->] (sx.south) -- (sx.south |- cat.north);
\draw[->] (sr.south) -- (sr.south |- cat.north);
\draw[->] (cat.south -| emb.north) -- (emb);
\draw[->] (cat.south -| fill.north) -- (fill);
\draw[->] (st) -- (lm);
% matching
\node[mt, text width=100mm] (S) at (2.6,-8.1) {$\mat{S}=\mat{Z}\mat{E}\T/\tau$ \texttt{(8192, 378)};\quad own entry $\to-\infty$;\quad
pool by class (logsumexp) $\to$ \texttt{(8192, 8)}: columns 0--6 man-made classes, column 7 natural};
\draw[->] (emb.south) -- (emb.south |- S.north);
\draw[->] (lm.south) -- (lm.south |- S.north);
\node[mt, text width=46mm, below left=6mm and -50mm of S] (l2) {level 2: $\softmax$ of columns 0--6\\ \texttt{(8192, 7)}};
\node[mt, text width=50mm, right=4mm of l2] (l1) {level 1: $\operatorname{LSE}(\text{0--6})-\text{col 7}+b$\\ $\ell$ \texttt{(8192,)},\ $P(M)=\sigma(\ell)$};
\draw[->] (S.south -| l2.north) -- (l2);
\draw[->] (S.south -| l1.north) -- (l1);
\node[ls, text width=46mm, below=4mm of l2] (ce) {CE, man-made rows ($\approx$4096)};
\node[ls, text width=50mm, below=4mm of l1] (bce) {BCE, all 8192 rows};
\draw[->] (l2) -- (ce); \draw[->] (l1) -- (bce);
\node[ls, text width=40mm, left=4mm of ce] (fl) {$0.5\,|\hat a-a|$, all rows};
\draw[->] (fill.west) -| (fl.north);
\end{tikzpicture}
\caption{Tensor shapes through one training step for one sensor, with \textbf{level~1 by matching} (the route through
column 7 of the pooled scores). Grey: the generator; teal: the shared stem, run three times; blue: the pixel tower;
green: the library tower; plum: matching; amber: losses. The dashed grey head is N001's linear level~1, still built,
no longer read. $B=16$, $P=8192$, $N=200$ \tagA, $L=378$.}
\label{fig:namer:shapes}
\end{figure}

\begin{worked}[label={we:namer:pool}]{From one row of $\mat{S}$ to both levels}
Shrink the library to six entries, $\tau$ folded into the scores \tagA: $e_1$ plastic 4.0 (the planted one), $e_2$
plastic 3.0, $e_3$ metal 2.0, $e_4$ concrete 1.0 (test split), $e_5$ soil 2.5, $e_6$ water 0.0. The sample is man-made,
class plastic.
\begin{align*}
\text{row 18 (exclude }e_1) &: (-\infty,\ 3.0,\ 2.0,\ 1.0,\ 2.5,\ 0.0)\\
\text{row 19 (classes)} &: (\text{plastic},\ \text{plastic},\ \text{metal},\ -1,\ \text{nat},\ \text{nat})\\
S_{\text{plastic}} &= \ln(e^{-\infty}+e^{3}) = 3.0,\qquad S_{\text{metal}}=2.0,\qquad S_{\text{concrete}}=-\infty\\
S_{\text{nat}} &= \ln(e^{2.5}+e^{0}) = \ln(12.1825+1) = \ln13.1825 = 2.5789\\
\ell &= \ln(e^{3}+e^{2})-2.5789 = \ln(20.0855+7.3891)-2.5789\\
&= 3.3133-2.5789 = 0.7344\\
P(M) &= \sigma(0.7344) = 1/(1+e^{-0.7344}) = 1/(1+0.4798) = 0.6758\\
P(\text{plastic}\mid M) &= e^{3}/(e^{3}+e^{2}) = 20.0855/27.4746 = 0.7311
\end{align*}
The test-split concrete entry $e_4$ is encoded but belongs to no class, so it plays no part; a class with no entry in
play scores $-\infty$ and gets probability 0.
\end{worked}

%======================================================================
\section{The loss, the optimiser and the schedule}\label{sec:namer:loss}

\begin{equation}
\mathcal{L}=\underbrace{\operatorname{BCE}(\ell,M)}_{\text{every sample}}
+\underbrace{\operatorname{CE}(S_{0..6},c)}_{\text{man-made samples}}
+0.5\,\underbrace{|\hat a-a|}_{\text{every sample}}
\label{eq:namer:loss}
\end{equation}
(\file{namer/train.py}, \code{loss}; \code{fill_weight} 0.5). \emph{BCE}, binary cross-entropy, is the loss for a
yes/no answer, $-[y\ln p+(1-y)\ln(1-p)]$; CE is minus the log of the right class's probability. Under matching $\ell$ is
\cref{eq:namer:l1}; under the old head it was \code{to_manmade}'s output. The fill term is on \emph{every} sample,
natural mixes and plain pixels too (plain target 0), not on man-made samples only as the proposal said; its ``As built''
list records the change. A fill head that learns ``nothing was added'' is level-1 evidence.

Two chance levels anchor any loss curve. A level-1 answer that knows nothing outputs $p=\frac12$:
$\operatorname{BCE}=-\ln\frac12=\ln2=0.6931$. A class answer that knows nothing over 7 classes: $\ln7=1.9459$. The
untrained namer at the sizes of \cref{sec:namer:step} gives BCE 0.7344 and CE 1.9941 on random data \tagC, close to both.
The CPU smoke run of the first code (EnMAP and SpectralEarth smoke shards, 120 steps, linear level~1) moved the class
CE from 1.81 to 1.02, level-2 accuracy at fill 1.0 from 0.62 to 0.65 on EnMAP, and left BCE ``still $\sim$0.68 (too
short to judge)'' \tagM{research/experiments/spectral-namer/N001/README.md}: level~1 had not left chance, an early sign
of what epoch~1 then confirmed.

\paragraph{Per step.} Each sensor's loss is multiplied by its weight (1) and divided by the number of sensors (3)
before \code{backward()}, so gradients from the three sensors add up within one step. Then the gradient norm is clipped
to 1, AdamW \citep{loshchilov2019decoupled} steps, and the scheduler steps. On a GPU the forward runs under bfloat16
autocast; the class scores are computed in float32 (\code{z.float()}, \code{e.float()}). A step sees up to
$3\times16\times512=24{,}576$ samples, about 491.5 million over $40\times500=20{,}000$ steps.

\paragraph{Two learning rates.} The stem and its BatchNorm (parameter names starting \code{stem.} or \code{norm.}) learn
at $0.3\times$ the base rate $3\times10^{-4}$, that is $9\times10^{-5}$; everything else, including $\ln\tau$ and $b$,
at $3\times10^{-4}$. The stem already knows spectra from V2; a slower rate keeps that knowledge while the new towers,
which start from random weights, catch up. Weight decay is 0.01 for both groups.

\paragraph{Warm-up and cosine.} The scheduler multiplies both rates by
\begin{equation}
\lambda(s)=\min\Big(1,\frac{s+1}{500}\Big)\times\frac12\Big(1+\cos\Big(\pi\min\big(1,\tfrac{s}{20000}\big)\Big)\Big),
\end{equation}
a linear warm-up over 500 steps times a half cosine over the whole run.

\begin{worked}[label={we:namer:lr}]{The schedule at a few steps}
\begin{align*}
s=0 &: \min(1,1/500)=0.002,\ \cos0=1,\ \lambda=0.002\times1=0.002,\ \text{lr}=6.0\times10^{-7},\\
s=249 &: 250/500=0.5,\ \tfrac12(1+\cos(0.0391))=0.99962,\ \lambda=0.49981,\ \text{lr}=1.499\times10^{-4},\\
s=10000 &: 1\times\tfrac12(1+\cos(\pi/2))=\tfrac12(1+0)=0.5,\ \text{lr}=1.5\times10^{-4},\ \text{stem }4.5\times10^{-5},\\
s=15000 &: \tfrac12(1+\cos(3\pi/4))=\tfrac12(1-0.70711)=0.14645,\ \text{lr}=4.39\times10^{-5}.
\end{align*}
The script prints the same \tagC.
\end{worked}

\paragraph{One BatchNorm, nine kinds of batch: why it is frozen.} In the first version of the code the post-stem
BatchNorm was in training mode during training. Each sensor's \code{forward} encodes the pixels, the rings and the whole
library in three separate calls, so it saw nine kinds of batch per step, each normalised by its own statistics. Its
stored running mean is an exponential average with momentum 0.1, $\text{running}\leftarrow0.9\,\text{running}+0.1\,
\text{batch}$; in steady state the call $k$ places from the end of a step carries weight $0.1\times0.9^k/(1-0.9^9)$,
and $1-0.9^9=1-0.3874=0.6126$. The last call, AVIRIS-NG's library, weighs $0.1/0.6126=0.1632$; by kind, pixels 0.299,
rings 0.332, library 0.369 (\file{research/textbook/tools/compute_namer_submission_output.txt}) \tagC. Evaluation and
inference would then normalise all three by that one blend: \cref{ch:finetune}'s BatchNorm see-saw in a new form, a
train/evaluation mismatch found by reading the code and never measured. \textbf{It was fixed before N001 ran} (code
bundle \texttt{namer2}): the norm never leaves evaluation mode; its running statistics stay V2's and its scale and shift
(128 numbers) still train, in the stem's group.

%======================================================================
\section{The library, the classes and the splits}\label{sec:namer:splits}

\paragraph{The bank.} \code{build_bank} (\file{namer/bank.py}) resamples the family library to one band list with the
library's own resampler (Gaussian spectral responses of width $1.1\times$ the band spacing), fills bands a spectrum does
not cover by linear interpolation along wavelength (and the nearest end beyond its range), keeps entries with at least
90\% finite bands, maps each to a class, and keeps the first entry of every repeated name. The result is a
\code{LibraryBank}: \code{refl} \shape{(L, N)}, \code{cls} \shape{(L,)} with $-1$ for natural, \code{test}
\shape{(L,)}, the names, and the held-out class.

\paragraph{The classes.} \file{namer/classes.py}: natural chapters $\to$ natural; building and pavement $\to$
concrete/asphalt; roofing dissolved by name (fiberglass, rubber, clear roof $\to$ plastic; cedar, shake, wood $\to$
wood/paper; metal, tin, steel, aluminium $\to$ metal; tar, felt, asphalt, shingle, tile $\to$ concrete/asphalt; any
other roofing $\to$ concrete/asphalt); glass left out (one spectrum, ``cannot be tested''). No roofing entry of the
shipped library maps to metal: it holds no metal roofing.

\paragraph{Held-out spectra.} A man-made entry is in the test split if the first 8 hexadecimal digits of the SHA-1 hash
of its name, read as a number and divided by $2^{32}-1=4{,}294{,}967{,}295$, fall below 0.2 (\code{_is_test}). Test
spectra are never planted in training and are not candidates in the class scores, in training or in evaluation. A hash
rather than a random draw means every band list agrees on the split without sharing a file.

\begin{worked}[label={we:namer:hash}]{Two metals, two sides}
\begin{align*}
\text{``Copper Metal''}: &\quad \texttt{1600007b} = 369{,}098{,}875,\\
&\quad 369{,}098{,}875/4{,}294{,}967{,}295 = 0.0859 < 0.2\ \Rightarrow\ \text{test};\\
\text{``Aluminum Metal''}: &\quad \texttt{d13908a4} = 3{,}510{,}175{,}908,\\
&\quad 3{,}510{,}175{,}908/4{,}294{,}967{,}295 = 0.8173\ \Rightarrow\ \text{train}.
\end{align*}
\end{worked}

\paragraph{De-duplication.} Some library names repeat (``Copper Metal'' three times, ``Galvanized Steel Metal''
twice). The hash sends repeats to one side, so the test split stays clean, but in training, leaving one ``Galvanized
Steel Metal'' out left its namesake in the class score, a near-copy that defeats the leave-out. Since the N001 code
bundle (\texttt{namer2}), \code{build_bank} keeps only the first entry per name.

\begin{table}[htbp]
\centering\small
\begin{tabular}{lrrrr}
\toprule
class & entries & held out & share & train\\
\midrule
oil/fuel & 40 & 8 & 20\% & 32\\
plastic & 46 & 8 & 17\% & 38\\
textile/fabric & 14 & 3 & 21\% & 11\\
wood/paper & 11 & 2 & 18\% & 9\\
paint/coatings & 4 & 1 & 25\% & 3\\
metal & 10 & 3 & 30\% & 7\\
concrete/asphalt & 34 & 9 & 26\% & 25\\
\midrule
all man-made & 159 & 34 & 21\% & 125\\
natural & 219 & -- & -- & 219\\
\midrule
bank & 378 & & & \\
\bottomrule
\end{tabular}
\caption{The bank on EnMAP's 224-band list, one entry per name (AVIRIS-NG: 380 entries, 220 natural, 160 man-made,
one more wood/paper entry, 34 held out) \tagC. Under level~1 by matching the 219 natural entries are class 7 in
training.}
\label{tab:namer:split}
\end{table}

\begin{sourcenote}[: about 20\% overall, not per class]
The hash ignores the class, so the held-out share varies: 17\% to 30\% per class with one entry per name, as the
docstring of \file{namer/bank.py} now says. The N001 README still says ``$\sim$20\% per class'', and its counts (408,
174, 41) are from before the de-duplication; the code gives 378, 159, 34 on EnMAP.
\end{sourcenote}

\paragraph{A held-out class (N002): textile/fabric.} \file{configs/namer/N002.json} holds the whole textile/fabric class
out (\code{"holdout_class": "textile_fabric"}). It is never planted in training and never a candidate; at evaluation it
is planted, and the trainer records how often level~1 calls it man-made and how often level~2 gives it a known class
with $P>0.5$. Pass: man-made at least 70\%, confident wrong at most 20\%. On EnMAP's bank this leaves
$125-11=114$ man-made training entries, and the evaluation pool of unseen man-made spectra is the 34 held-out spectra
plus textile's 11 training-split ones, 45 \tagC.

\begin{sourcenote}[: textile, not paint]
N002 first held out paint/coatings. After de-duplication paint has 4 spectra, too few to test anything, and on
2026-10-31 the user chose textile/fabric (14 spectra) instead; \file{configs/namer/N002.json} says so. Three places
still say paint: the proposal's ``As built'' list (``Open-set class: paint/coatings (run N002)''), its ``Splits'' text
(``e.g. paint/coatings''), and the copy \file{research/experiments/spectral-namer/N002/config.json}, which is the old
config. Likewise \file{configs/namer/N001.json} now holds N001b's settings (run directory \texttt{N001b},
\texttt{"level1": "matching"}); N001's own config survives as
\file{research/experiments/spectral-namer/N001/config.json}.
\end{sourcenote}

%======================================================================
\section{The test it must pass}\label{sec:namer:pass}

Each epoch \code{evaluate} plants test-split spectra into test shards at fills 0.1, 0.25, 0.5 and 1.0, half man-made
and half natural mix (\code{kinds=(0.5, 0.5, 0.0)}), 40 batches with a fixed seed: up to $40\times8192=327{,}680$
samples per sensor and fill \tagC. Level~1 can buy recall by calling more things man-made, so the fair comparison holds
false alarms equal: read ACE's natural false-alarm rate $f$ from \file{separability_<sensor>.json}
($1-$\code{level1_natural_called_natural}), cut the namer's logits at the $(1-f)$ quantile of its natural plantings, and
count the man-made plantings above the cut:
\begin{lstlisting}
cut = torch.quantile(torch.cat(nat_logits), 1 - fa)
recall = (torch.cat(mm_logits) > cut).float().mean()
\end{lstlisting}

\begin{worked}[label={we:namer:quantile}]{Recall at a fixed false-alarm rate}
Ten natural logits, sorted: $-3,-2.5,-2,-1.5,-1,-0.5,0,0.5,1,2$; $f=0.2$ \tagA. The 0.8 quantile, with linear
interpolation as in \cref{sec:material:mask}, sits at position $0.8\times(10-1)=7.2$, between 0.5 and 1.0:
\begin{align*}
\text{cut} &= 0.5+0.2\times(1.0-0.5) = 0.5+0.1 = 0.6 .
\end{align*}
Five man-made logits $-1, 0.4, 0.7, 1.5, 3$: above 0.6 are three, so recall $=3/5=0.6$ \tagC.
\end{worked}

\begin{table}[htbp]
\centering\small
\begin{tabular}{lcccc}
\toprule
sensor & ACE natural right & $f$ & ACE man-made recall & namer needs\\
\midrule
PRISMA & 0.8342 & 0.1658 & 0.7724 & $\ge0.8724$\\
EnMAP & 0.7851 & 0.2149 & 0.6635 & $\ge0.7635$\\
AVIRIS-NG & 0.8134 & 0.1866 & 0.7603 & $\ge0.8603$\\
\bottomrule
\end{tabular}
\caption{Level~1's bars at fill 0.25: ACE $+0.10$ at ACE's own false-alarm rate, from the three separability JSONs
\tagC.}
\label{tab:namer:bars}
\end{table}

Level~2 must also beat ACE by 0.10, but not against these JSONs yet: their level-2 numbers were measured on the 17 old
families and ``must be re-scored on the v1 classes \dots\ before comparing; that re-score is the first job after
launch'' \tagM{research/experiments/spectral-namer/N001/README.md}. \file{best.pt} is the epoch with the best mean
level-2 accuracy at fill 0.25 over the sensors, and every epoch also writes \file{ep<n>.pt} and a row of
\file{metrics.csv} (losses, $\tau$, learning rate, and per sensor and fill the level-1 and level-2 accuracies, the
recall at ACE's rate, and for N002 the two open-set rates).

\begin{sourcenote}[: three things the proposal lists and the trainer does not do]
It computes no bootstrap interval (the pass rule's ``interval above zero'' is a later step); it has no ``unknown''
output (level~2 always returns its best class and \code{p_cls}; ``class unknown'' needs a threshold not yet chosen);
and natural spectra are planted in both splits, so level-1 natural accuracy is not an unseen-spectrum test (``As
built'').
\end{sourcenote}

%======================================================================
\section{Inference: naming the pixels of a mask}\label{sec:namer:infer}

\code{load_namer} (\file{namer/infer.py}) rebuilds the model from a checkpoint (\file{last.pt}, \file{best.pt} or
\file{ep<n>.pt}): the normalisation curve from the stored \code{stats.*} buffers, the sizes from the stored config's
\code{model} entry, and the level-1 mode from its \code{level1} entry (default \code{"head"}). Then
\begin{lstlisting}
name_pixels(model, cube, wavelengths, pixels, library_root, cache_dir, valid)
\end{lstlisting}
names the pixels of a boolean mask:
\begin{enumerate}
\item \textbf{Bands.} Keep a band if it is valid at more than half of the requested pixels
(\code{valid[:, rows, cols].mean(axis=1) > 0.5}): 3 valid of 5 pixels is 0.6, kept; 2 of 4 is exactly 0.5, dropped.
\item \textbf{Pixels and rings.} $\vect{x}$ is the cube at the mask's pixels, \shape{(P, N)}; the ring means are
computed one band at a time over the whole scene (\code{ring_means}, box sums of $15\times15$ minus $5\times5$ with
\code{uniform_filter}), so a full scene never needs a second cube in memory. The per-pixel band validity is passed to
the stem, which masks dead bands exactly.
\item \textbf{Library.} \code{build_bank} at the scene's bands, the \emph{whole} library, test-split spectra included:
at inference every spectrum is fair to match against. One \code{forward_library} call encodes it.
\item \textbf{Scores}, in batches of 4,096 pixels: \code{forward_pixel}, then \code{class_scores} with $C+1=8$ classes
under matching (all man-made entries in their classes, natural as class 7), \code{level1_logit} for level~1 and the
first seven columns for level~2; with the old head, 7 classes and the head's logit.
\item \textbf{Outputs} per pixel: $P(M)=\sigma(\ell)$; the best class and its $P(k\mid M)$ from the softmax of the seven
class scores; $\hat a$.
\end{enumerate}
No entry is excluded at inference: the pixel's own spectrum is unknown, so there is nothing to leave out.

\paragraph{Into the AIGC table.} \file{scripts/aigc_object_table.py} is the first caller (\cref{ch:submission}). With
\code{--namer-ckpt} it loads the vendable, passes its \code{normalized_hyperspectral_cube} (reflectance, despite the
name; \file{docs/09-known-issues.md}), \code{band_cw_order} and validity cube to \code{name_pixels} with the committed
anomaly mask, and then averages per 8-connected object: mean $P(M)$, the majority class and its share, mean
$P(k\mid M)$, mean fill. The findings file labels the namer column ``preliminary'', and objects with mean $P(M)>0.5$
count as man-made when it is compared with the family match.

\begin{keypoint}
Leave the planted spectrum out of its own class score, and a two-tower network can only lower its loss by placing a
whole material together; read \emph{both} questions from the same cosine table, so that level~1 compares the man-made
shelf with the natural one; and judge level~1 at the classical method's own false-alarm rate, so that recall cannot be
bought.
\end{keypoint}

\begin{drift}
Everything in this chapter: N001b and N002's settings (\file{configs/namer/}), the bank counts (378 / 219 / 159 / 34 on
EnMAP), the per-class shares, the 347,685 parameters, the two early level-1 measurements, the CPU smoke numbers, and the
pass bars. N001b's and N002's results were pending when the chapter was written (2026-10-31).
\end{drift}

%======================================================================
\begin{confusions}
\confusion{The namer reuses Indradhanu~V2, so it detects anomalies.}{It reuses V2's stem, 5\% of its parameters, as a
spectral feature extractor. It names pixels something else has already flagged.}
\confusion{Siamese means the two sides are identical.}{Only the stem is shared, 17,825 weights. After it the pixel tower
($214{,}402$ weights, ring input, fill head) and the library tower ($115{,}328$) differ on purpose.}
\confusion{Logsumexp is the best cosine.}{A smooth maximum: $\ln(e^7+e^6)=7.31$, not 7. Every member counts, the
closest most.}
\confusion{Level~1 by matching needs a new model.}{It needs one new number, $b$, and a different reading of the same
score table. N001's head is still in the model, unused.}
\confusion{$P(M)$ is the best man-made cosine against the best natural one.}{It is a share of softmax mass
(\cref{eq:namer:l1share}): every entry contributes, weighted by $e^{s_j}$, and the side with more entries starts ahead by
$\ln$ of the count ratio, which $b$ can absorb.}
\confusion{Leaving the planted spectrum out is an evaluation trick.}{Here it is a training rule too, for man-made
\emph{and} natural plantings, and the main departure from the proposal's per-spectrum InfoNCE.}
\confusion{The ring is 2 to 7 pixels away.}{Beyond 2: Chebyshev distance 3 to 7, 200 pixels in the interior.}
\confusion{20\% of every class is held out.}{About 21\% of the man-made spectra overall on EnMAP (34 of 159); 17\% to
30\% per class.}
\confusion{N002 holds out paint.}{It did in the first config. It holds out textile/fabric (14 spectra) since the user's
decision of 2026-10-31.}
\confusion{A frozen BatchNorm does not learn.}{Its running statistics are frozen; its scale and shift still train.}
\end{confusions}

\begin{checkyourself}
\item\label{ex:namer:1} In \cref{we:namer:ring}, make the halo 0.60 and the target 0.90. What is the ring mean, and why?
\item\label{ex:namer:2} With $a_{\min}=0.02$ instead of 0.05, what share of fills is at most 0.25, and what is the mean
  fill? Use \cref{eq:namer:loguniform} and the mean formula.
\item\label{ex:namer:3} Plant $\vect{t}=(0.30,0.10)$ into $\vect{b}=(0.20,0.40)$ at $a=0.5$ with $\vect{g}=(1.1,0.9)$
  and no noise.
\item\label{ex:namer:4} In \cref{we:namer:classscore}, suppose metal had only $e_1$ and $e_2$. With $e_1$ left out,
  compute $S_{\text{metal}}$, $P(\text{metal})$ and the loss.
\item\label{ex:namer:5} Count the parameters of a namer with \code{hidden} 128 and \code{embed} 64, everything else as
  N001b. (Hint: only the towers change.)
\item\label{ex:namer:6} At PRISMA fill 0.25 the namer, cut at $f=0.1658$, recalls 85\% of man-made plantings. Does
  level~1 pass? What recall would it need?
\item\label{ex:namer:7} Show that the level-2 gradient entries $q_j-\mathbf{1}[j\in c]w_j$ always sum to zero.
\item\label{ex:namer:8} In \cref{we:namer:l1}, add a second natural entry, water, with $s=5.5$. Compute $\ell$, $P(M)$,
  the BCE for a man-made truth, and the four gradients. Which side gained from the extra entry?
\item\label{ex:namer:9} Bands at 1790, 1800, 1960, 1970\,nm. Compute their widths, and the bins they fall in.
\item\label{ex:namer:10} Level~1 by matching with $b=0$ on a bank of 219 natural and 125 man-made entries: a pixel has
  the same score $s$ against every entry. What is $P(M)$? What $b$ would make it $\frac12$?
\end{checkyourself}

\nextchapter{ch:submission}
